% appendix_e_pilots.tex -- Appendix E Pilot results in full. Owner: Appendix writer.
% Sources (V1, 100-image AnnyOne holdout, matched people, unless noted):
%   depth pilots  = exp1: docs/handover.md 3.1 (table and "What this says"); success rule docs/experiments_plan.md:93-95
%                   (71.5 mm, 89.2 %); full table = tables/T08_depth_runs.tex (tab:depth-runs), not repeated here.
%   sampler pilot = exp2: H 3.4 "The sampler (exp2)"; rule docs/experiments_plan.md:135-139 (child PVE limit = \DiagChildSwapBefore);
%                   H 6 "Children".
%   head-loss pilot = exp3: H 3.2 "Refuted" (3571957:98, 3571960:318, 3571958:99; neck group 62.5, 3571957:100);
%                   rule docs/experiments_plan.md:179-180; safety limit missed 71.8 mm / 89.0 % (deck explanation slide 22).
%   finger pilot  = exp4: H 3.3 (3571959:316, 3571960:287-314); rule docs/experiments_plan.md:219-221.
%   geodesic pilot = exp5: H 3.2 "Refuted" (3573521:89-91, :98, :103); rule docs/experiments_plan.md:404-410
%                   ("flagged" = bones the Base dump labels FLIP_*, PARTIAL_FLIP or CONST_OFFSET).
%   bone-correction pilot = exp6: H 3.2 tables (3579014:12, :37, :101-121; 3579015:70-86); rule
%                   slurm/experiments/exp6_bone_fix.sh:48-52; docs/experiments_plan.md:524-528.
%   noise: tables/T14_pilot_noise.tex (tab:pilot-noise), H 3.6; deck explanation slide 22.
%   history: T20 header; stages T19. Step times (\MethodItTime...) from the draft, no kept log; moved from 04 (6 Nov).
% Numbers come from numbers_diagnostics.tex (\Diag...) where a macro exists; the rest are written out with the source
% in the comment of the same line. No job ids, no old run names (plan section 5).
% Trimmed 19 Oct: the six per-pilot paragraphs after Table E.1 dropped (Appendix F holds the diagnosis); recover with
% git show 76aa1a4:docs/thesis/chapters/appendix_e_pilots.tex.
% 6 Nov: the two "lessons" passages merged into one section with both labels.
%   terms: age bias = mean(pred - GT) of the age phenotype, scripts/diagnostics/diag_report.py:206,314;
%   flagged = docs/experiments_plan.md:404-405; geodesic = 05_diagnostics sec:diag-neck.
% 7 Nov: audit fixes. The draft's "six times more images" had no source (and T19's sample counts do not give it):
%   replaced by the Student's data (A+P, v23 labels; T19). Typed 71.5 / 89.2 may become \DiagRulePVE /
%   \DiagRuleRecall at integration (not in BASE).

\chapter{Pilot results in full}
\label{app:pilots}

Before the final runs I tested six changes in pilots. A pilot is a short run, usually a 2-epoch fine-tune of Base,
that tests one change. I compare each pilot with a control, which is the same training without the change, and
Table~\ref{tab:pilots-all} names the control of each pilot. Appendix~\ref{ch:diagnostics} explains what each result
means. The second section is older history, the ViT-S students I trained first and the stages that produced Base.

\section{The six pilots}

Every pilot fine-tuned Base on AnnyOne only, at learning rate $10^{-5}$ and batch 4, for 2 epochs (8,000 images)
unless noted. I evaluated all of them on the 100-image AnnyOne holdout in setting V1. Table~\ref{tab:pilots-all} gives the success
rule that I wrote into each script before the run. The full depth runs are in Table~\ref{tab:depth-runs}.

The table uses three terms that have not come up before. The geodesic angle is the angle of the single rotation that
turns the predicted bone into the true one. The loss I normally train with does not use this angle. It compares the
two rotation matrices element by element with L1. A flagged bone is one that the per-bone dump of Base in
Appendix~\ref{ch:diagnostics} marks as flipped or as turned by a constant offset. Age bias is the mean of the
predicted minus the true age phenotype, so it says whether the predicted age is too high or too low on average.

\begingroup
\footnotesize
\setlength{\tabcolsep}{3pt}
\begin{longtable}{@{}>{\raggedright}p{1.6cm}>{\raggedright}p{2.6cm}>{\raggedright}p{2.7cm}>{\raggedright}p{2.95cm}>{\raggedright\arraybackslash}p{2.1cm}@{}}
\caption[Pilot experiments: change, rule, result, control]{The six pilot experiments. V1, 100-image AnnyOne holdout,
matched people; fine-tunes of Base, one run each, without a seed. The rule is the success rule written into each
script before the run. Head direction is the global head orientation error.}
\label{tab:pilots-all}\\
\toprule
Pilot & What changed & Success rule & Result & Control \\
\midrule
\endfirsthead
\toprule
Pilot & What changed & Success rule & Result & Control \\
\midrule
\endhead
\bottomrule
\endfoot
Depth pilots (2 and 10 epochs) &
training targets at the head joint instead of the pelvis; depth and translation weights $\times$1, $\times$5, $\times$20 &
slope toward 1 and intercept toward 0\,m, with PVE $\le$ 71.5\,mm and recall $\ge$ 89.2\,\% &
rule not met (PVE or recall missed in every run, controls too); root error \DiagRootErrBase{} $\to$ \DiagXFiveTenRoot\,m, size \DiagSizeBase{} $\to$ 0.84, slope \DiagDepthSlopeBase{} $\to$ \DiagSlopeStrongLow{} ($\times$5, 10 epochs) &
same schedule, targets at the pelvis, weight $\times$1 (one 2-epoch, one 10-epoch) \\
\addlinespace
Sampler pilot &
age-balanced sampling: children (age $<$ 0.3) per batch \DiagSamplerBefore{} $\to$ \DiagSamplerAfterLow--\DiagSamplerAfterHigh\,\% &
10 manifest children: PVE $<$ \DiagChildSwapBefore\,mm and $|$age bias$|$ $<$ 0.078; holdout PVE $\le$ 71.5, recall $\ge$ 89.2 &
inconclusive: children \DiagSamplerChildPVE{} vs.\ \DiagSamplerChildPVECtrl\,mm, age bias 0.078 vs.\ 0.085; baby age floor \DiagSamplerBabyBefore{} $\to$ \DiagSamplerBabyAfter &
same run with the sampler off (age $\times$3 and height $\times$2 shape weights in both) \\
\addlinespace
Head-loss pilot &
L1 loss on the camera-frame orientation of head, neck03 and both wrists, weight 5 &
head direction $<$ 12.0\,deg; PVE $\le$ 71.5\,mm, recall $\ge$ 89.2\,\% &
head direction \DiagGlobalHeadArm\,deg, within the controls' \DiagGlobalHeadCtrlA--\DiagGlobalHeadCtrlB; safety limit missed (71.8\,mm, 89.0\,\%) &
two identical runs at weight 0 \\
\addlinespace
Finger pilot &
region-normalised rotation loss (body, hands, head-neck averaged separately); a test-time person crop &
fingers below the rest-pose floor \DiagFingersRest\,deg without a PVE cost &
no effect: fingers \DiagFingersRegion{} vs.\ \DiagFingersRegionCtrl\,deg; crop \DiagFingersZoom\,deg at body PVE \DiagZoomPVE\,mm &
same run with the loss off (one of the two identical runs above) \\
\addlinespace
Geodesic pilot &
geodesic angle instead of L1 on the rotation matrices &
flagged bones down by $\ge$ 3\,deg more than the control, or neck and head under 15\,deg &
no movement: neck02 \DiagNeckTwoBase{} $\to$ \DiagGeodesicNeckTwo, head \DiagDumpHead{} $\to$ \DiagGeodesicHead\,deg &
same run with the L1 loss \\
\addlinespace
Bone-correction pilot &
constant output rotation on neck02, head, metacarpal1.L, finger4-1.R; no training, then a 2-epoch fine-tune &
bones $<$ 15\,deg, girth $\ge$ 0.9; no training: head direction $\le$ Base + 3\,deg; fine-tuned: guards against the control &
no training: fails (head direction \DiagHeadDirNoTrain{} vs.\ limit 15.0\,deg); fine-tuned: passes (neck02 4.7\,deg, girth 0.997, fingers \DiagFingersFixFT\,deg) &
the geodesic pilot's L1 control \\
\end{longtable}
\endgroup

\paragraph{Pilot variation.} Two runs of the same short configuration differed by \DiagNoisePVE\,mm PVE,
\DiagNoiseHead\,deg head direction and \DiagNoiseRecall{} recall points
(Table~\ref{tab:pilot-noise}). The restart comparison also showed higher average PVE than Base.
These observations warn against reading too much into small single-run differences. They do not
give a universal significance threshold. I selected the placement and bone corrections using their
targeted diagnostics and matched controls, rather than a small PVE gain.

\section{The ViT-S students and the stages that produced Base}
\label{sec:method-kd-history}
\label{sec:res-history}

% T20 History: the staged schedule to Base and the ViT-S students (Results 6.11). Owner: Results writer.
% Source: the draft's tab:main (draft main.tex:674-694; tables/draft_main.tex) and tab:kd_eval (main.tex:789-810;
% tables/draft_kd_eval.tex). All numbers are setting V0 (helper bones unmasked) on the 100-image AnnyOne holdout,
% measured before the depth fix, the bone correction and the corrected 3DPW labels.
% Dropped on purpose: (1) the draft's TODO cells of tab:main (v2-v4 never re-evaluated, plan 7 Q6) -> "--";
% (2) the v10 row of tab:kd_eval (87.4 mm V0 at epoch 299 cannot be reconciled with the V1 74.1 mm at epoch 290
% of H 3.1 without the logs; v10 is not a named model of this thesis); (3) the 3DPW table tab:test3dpw (old labels,
% retired). Its v23 re-evaluation of S6c and S8 (job 3615020, 3615020_2:119 and 3615020_4:118, the script of
% tab:final-3dpw) is the two 3DPW columns of part (b); the same jobs re-measured the AnnyOne cells of S6c and S8 and
% gave exactly the draft's values (114.5 / 84.6 / 118.4 / 81 / 89; 120.6 / 86.9 / 125.0 / 76 / 87). Base and Mixed-L
% 3DPW cells: tab:final-3dpw (3612073, 3612074).
% Cross-checks: Base 74.3 / 57.5 / recall 90 and Mixed-L 78.5 / 57.4 / 89 equal the AnnyOne test-100
% V0 numbers of the v23 3DPW-test logs (3612073:25-32, 3612074:25-32; H 0b).
% Names: stages as in Chapter 4 (stage 4 = Base); students S1-S8 = v1-v8 in the code (numbers_method.tex).
% 12 Nov 2026: V0 -> unmasked protocol; captions to at most three sentences, variant descriptions moved to a table
% note; Base row labelled as the initial model.
\begin{table}[!htbp]
\centering\footnotesize
\setlength{\tabcolsep}{2pt}
\begin{subtable}{\linewidth}
\centering
\begin{tabular}{@{}llrrrr@{}}
\toprule
Run & Change against the previous run & PVE & PA-PVE & MPJPE & Rec. \\
\midrule
L-stage 1 & fresh heads, all blocks trained (side branch) & -- & -- & -- & -- \\
L-stage 2 & fresh heads, backbone frozen & \HistStageOnePVE & -- & -- & -- \\
L-stage 3 & + last 4 of 24 blocks, helper mask & \HistStageTwoPVE & -- & -- & -- \\
L-stage 4 & + neck and hand loss $\times$4 & 73.6 & 56.5 & 73.2 & -- \\
Base & + shape loss weight 20 & \BaseVzeroPVE & 57.5 & 74.4 & 90 \\
\bottomrule
\end{tabular}
\caption{The transfer schedule that produced the initial model, Base (Table~\ref{tab:model-names}); recall (Rec.)
in \%. L-stages 2--4 are evaluated during training, the initial model with the standalone script.}
\label{tab:history-stages}
\end{subtable}

\vspace{0.8em}
\begin{subtable}{\linewidth}
\centering
\begin{tabular}{@{}llrrrrrrrr@{}}
\toprule
 & & & \multicolumn{5}{c}{AnnyOne holdout, unmasked} & \multicolumn{2}{c@{}}{3DPW test, v23$^{c}$} \\
\cmidrule(lr){4-8}\cmidrule(l){9-10}
Model & Teacher & Params & PVE & PA-PVE & MPJPE & Prec. & Rec. & PVE & PA-PVE \\
\midrule
Base (ViT-L) & -- & 318.4 & \BaseVzeroPVE & 57.5 & 74.4 & 81 & 90 & \BaseThreeDPW & \BasePAThreeDPW \\
Mixed-L & -- & 318.4 & \BaseOldVzeroPVE & 57.4 & 78.6 & 83 & 89 & \BaseOldThreeDPW & \BaseOldPAThreeDPW \\
\addlinespace
Student-S/2 & Base & 32.8 & 193.4 & 128.2 & 210.7 & -- & 86 & -- & -- \\
Student-S/6b$^{a}$ & Base & 37 & 126.9 & 92.1 & 132.0 & 77 & 88 & -- & -- \\
Student-S/6c & Base & 37 & \textbf{114.5} & \textbf{84.6} & \textbf{118.4} & 81 & 89 & \StudentSSixcThreeDPW & \StudentSSixcPAThreeDPW \\
Student-S/mixed$^{b}$ & Base & 37 & 120.6 & 88.2 & 125.2 & 78 & 87 & -- & -- \\
Student-S/final$^{b}$ & Mixed-L & 37 & 120.6 & 86.9 & 125.0 & 76 & 87 & \StudentSFinalThreeDPW & \StudentSFinalPAThreeDPW \\
\bottomrule
\end{tabular}

\smallskip
\parbox{\linewidth}{\scriptsize Student-S/2: feature distillation only; Student-S/6b: heads fed through a projector at
the teacher's width; Student-S/6c: Student-S/6b at batch 16; Student-S/mixed and Student-S/final: continued with 3DPW
images. $^{a}$Evaluation during training. $^{b}$30\,\% 3DPW images per batch, old labels. $^{c}$Corrected labels,
the script of Table~\ref{tab:final-3dpw}; the teacher (Final) reaches \FinalThreeDPW{} / \FinalPAThreeDPW\,mm.}
\caption{The ViT-S students against their ViT-L teachers; parameters in millions, precision (Prec.) and recall
(Rec.) in \%.}
\label{tab:history-students}
\end{subtable}
\caption[History: the schedule to the initial model and the ViT-S students]{Earlier models, measured before the
depth fix, the bone correction and the corrected 3DPW labels, on the AnnyOne holdout (100 images) in the unmasked
protocol (helper bones scored as predicted); errors in mm, single runs without a seed. ``--'': not measured with the
same script; the students' old-label 3DPW numbers are retired.}
\label{tab:history}
\end{table}


I trained Base in stages (the names are in Table~\ref{tab:names}, the numbers in Table~\ref{tab:history-stages}).
L-stage 1 trained new Anny heads together with every block of NAVER's ViT-L backbone. I did not continue it, and no
checkpoint of it is kept. L-stage 2 started again from NAVER's backbone, kept it frozen and trained only new Anny
heads. From there each stage started from the one before. L-stage 3 unfroze the last 4 of the 24 blocks and added the
helper mask. L-stage 4 weighted the loss on the neck and hand bones four times. Finally, Base added the shape loss
weight of 20.

Before the ViT-B Student I tried smaller students to meet the size requirement. I trained a series of ViT-S students,
Student-S/1 to Student-S/final, with Base as their teacher. ViT-S is the vision transformer one size below the
Student's ViT-B. All runs in Table~\ref{tab:history}, the stages and the students, came before the depth fix, the bone
correction and the corrected 3DPW labels. So the table shows the numbers I measured at the time, in setting V0 (the
students are in Table~\ref{tab:history-students}). The best student, Student-S/6c, had \StudentSixcRatio{} times the
PVE of its teacher.

Four things I learned from these students went into the Student's training. With feature distillation alone, nothing
trains the heads, and Student-S/2 has by far the largest error in Table~\ref{tab:history-students}. For this reason
the Student starts with a copy of the teacher's heads. The copied heads only worked when the projector fed them their
input. With a new, randomly initialised input layer in front of them they did not. Heads at the teacher's width also
did better than narrower ones. So the Student keeps every head of Final at its full width. The starting point
mattered most. A backbone that NAVER had already trained for mesh recovery gave the largest gain, which is why
the Student starts from NAVER's ViT-B. Last, I increased the data budget and used corrected 3DPW labels alongside AnnyOne. The trials
do not prove that data, rather than capacity, was the only limitation.

Two settings made the ViT-S training steps faster. I ran the matrix products in TF32, a reduced-precision mode of the
A100, and I kept the teacher's backbone in bfloat16, a 16-bit number format. Together they cut a ViT-S student's
training step from \MethodItTimeBefore{} to \MethodItTimeAfter\,s on an A100, and the evaluation numbers did not
change. The Student's run keeps both settings: TF32 as in Final (Table~\ref{tab:flags-final}) and the teacher in
bfloat16 (Table~\ref{tab:flags-student}). I took these two step times from an earlier version of this report; the
training log behind them was not kept.
